\section{问题、两个例子与符号约定}
\emph{要判断系统能做什么，先看哪些信息没有提供给它。}
一个动力系统在什么条件下能恢复关于自身过去的一个事实？需要向它传递多少比特，又能通过什么干预检验它是否使用了这些信息？首先要明确，在提出查询的时刻，接收方可以获取哪些量。下面两个关于规则~110 的构造说明了这个问题的重要性。

\paragraph{两个传感器覆盖了整个窗口，却无法区分两段历史。}
将传感器置于 $S=(-1,1)$，在时刻 0 和 1 读取它们的值。任务是预测时刻 2 元胞 0 的值。传感器的依赖锥覆盖了该元胞的整个依赖窗口 $[-2,2]$。在第一种初始状态下，只有位置 $-1$ 处为 1。在第二种初始状态下，位置 $-1$ 和 $0$ 处均为 1。在这两种情况下，传感器的两次读数都是同一对值 $(1,0)$。时刻 2 的目标值却分别为 0 和 1。解码器依据这份记录作出的回答，必定在两种情况中的一种出错（第~\ref{sec:observer-coverage} 节）。

\paragraph{一个比特可以恢复已经消失的差异。}
现在，目标 $Y$ 是元胞 5 过去的取值。在两对历史中的每一对里，仅翻转 $Y$ 就会改变时刻 2 元胞 3 的读数 $Z$。观察者保存这一比特。到时刻 4，每对历史中的两个主格点系统构型已在整条直线上完全一致。把 $Z$ 写入元胞 $-1$。它的左邻居为 1，右邻居保存着上下文比特 $K$。下一次更新得到 $Z\oplus K=Y$。在格点系统状态相同的情况下，若改为送回 $1-Z$，四段历史的答案都会翻转（第~\ref{sec:observer-return} 节）。

在这一构造中，信道传送的是传感器的原始读数。在两步 OLMo 实验中，发送方先根据过去的选择 $h$ 和反馈 $e$ 计算修正后的规则 $p=h\oplus e$。进入新一轮的只有一个词元 $m\in\{\mathrm{A},\mathrm{B}\}$，它在呈现字符串 $\mathbf{x}$ 之前就已选定。接收方使用全新的 KV 缓存重新处理这个词元。因此，该实验测量的是一个已选定的规则能否跨轮次传递。要检验状态连续保留时系统能否在内部修正策略，还需另行实验（第~\ref{sec:olmo} 节）。

\begin{center}\small
\begin{tabularx}{\linewidth}{@{}lX@{}}\toprule
符号 & 全文含义 \\\midrule
$\omega,H,W,Z$ & 一段允许的历史、接收方掌握的信息、观察者可获取的信息和实际送达的消息。\\
$Y$ & 关于过去的目标比特。在 T1--T3 更一般的表述中，其取值范围可以是任意有限字母表。\\
$t,i,j,s$ & 离散时间与格点的整数坐标。$e_j$ 仅翻转初始元胞 $j$。\\
$B_t,\widetilde B_t,S_t$ & 未受扰动的构型、受扰动的构型，以及二者取值不同的坐标集合。\\
$G_t,D_t,I_t$ & 原点处元胞的值、本质初始坐标的集合，以及理论依赖窗口。\\
$h,e,p,m$ & OLMo 过去的选择、反馈比特、修正后的规则和消息词元。\\
$\mathbf{x}=LwR$ & 新的八比特输入。$w$ 是中间的六个比特，$L,R$ 是两端的比特。\\
$B,E(a,t)$ & 桥接公设组，以及实例 $a$ 在时刻 $t$ 是否具有现象体验的谓词。\\
\bottomrule\end{tabularx}
\end{center}

公式给出可由指定前提推出的结论（\texttt{COMPUTED}）。关于某一具体系统的结论需要可执行的见证或测量报告（\texttt{MEASURED}）。\texttt{UNBOUND} 表示尚未建立所需的证据关联，\texttt{NOT\_RUN} 表示在所述版本中尚未执行该项实验。下文的“精确”或针对给定数学定义域中的每个元素，或针对明确指定的有限测试集。第~\ref{sec:observer-theory} 节首先确定精确解码器存在的条件。
